\section{Le développement durable}\label{le-developpement-durable}

\stanceinfo{\textbf{Congrès de la FCEG 2018} [2018-01-07]}{2024-01-02}

\officialstance{La Fédération canadienne étudiante de génie estime que le développement durable est une considération essentielle dans la pratique de l'ingénierie et qu'il est nécessaire d'éduquer et d'impliquer ses membres sur les questions de développement durable, tout en évaluant et en améliorant le programme d'études sur le développement durable dans le domaine de l'ingénierie.}

\stanceheading{La position des étudiants}
\begin{itemize}
 \item Les ingénieurs ont la responsabilité de faire primer la sécurité, la santé et le bien-être du public et de tenir compte de l'environnement, ce qui implique la responsabilité d'effectuer un travail durable pour les générations futures.
 \item Les futures générations d'ingénieurs devraient recevoir une formation approfondie et s'investir dans les pratiques de développement durable de l'ingénierie avant l'obtention de leur diplôme, par le biais d'applications appropriées dans le cadre du programme d'études d'ingénierie.
 \item La FCEG devrait évaluer la durabilité de ses opérations afin de réformer et d'améliorer ses pratiques si nécessaire.
 \item Le travail de plaidoyer autour de la durabilité doit être évalué dans l'optique des étudiants en ingénierie, des opportunités d'emploi et de l'avenir de l'ingénierie.
\end{itemize}

\stanceheading{Les recommandations aux partenaires, aux parties intéressées et aux autres entités}
\begin{itemize}
 \item La FCEG incite le BCAPG à intégrer des unités d'accréditation axées sur le développement durable dans tous les diplômes d'ingénierie, et à garantir que l'ensemble des étudiants en ingénierie diplômés possèdent ces compétences.
 \item La FCEG encourage ses membres à plaider en faveur de l'importance du développement durable dans leurs activités quotidiennes au sein de leur propre institution, tout en les incitant à participer aux conférences respectives de la FCEG. Cela inclut la promotion de l'idée d'un désinvestissement institutionnel et la fourniture de ressources dédiées au développement durable.
 \item La FCEG encourage les DDIC à soutenir les initiatives, les équipes et les pratiques de développement durable au sein de leurs institutions respectives. Elle les incite également à établir des partenariats avec des organisations en vue de créer un campus d'ingénierie plus durable.
\end{itemize}

\stanceheading{Ce que fait la FCEG}
\begin{itemize}
 \item La FCEG organise chaque année une conférence nationale sur le développement durable en ingénierie dans l'une de ses écoles membres, dans le but d'éduquer et d'engager directement les étudiants en ingénierie sur les pratiques durables liées à la profession, qu'ils pourront rapporter et mettre en œuvre des idées dans leurs universités.
 \item La FCEG continuera à mettre à jour les meilleures pratiques en matière de développement durable pour son usage et celui de ses sociétés membres, y compris La Banque des meilleures pratiques en matière de développement durable, ainsi que d'autres ressources d'apprentissage utiles liées au développement durable.
\end{itemize}

\stanceheading{Ce que la FCEG envisage de faire}
\begin{itemize}
 \item La FCEG effectuera une analyse annuelle de l'empreinte environnementale de ses activités et réformera ses activités pour les aligner sur les pratiques durables dans la mesure du possible.
 \item La FCEG contribuera à l'élaboration d'une documentation sur les lignes directrices de la conférence en matière de développement durable, afin de prendre en compte les pratiques durables dans les processus de planification et d'exécution. En tant qu'organisation nationale d'étudiants, la FCEG a un rôle majeur à jouer pour ouvrir la voie à des actions durables en tant qu'organisation et pour motiver ensuite les écoles membres.
 \item La FCEG peut élaborer des documents d'orientation pour elle-même et ses organisations membres afin de réduire leur impact sur l'environnement, par le biais d'un document évolutif décrivant les meilleures pratiques. Ce document sera revu et mis à jour chaque année afin de refléter tout changement dans les meilleures pratiques.
\end{itemize}

\newpage
